\section{Experimental methodology} 
\label{sec:experiment} 
We decouple the contribution of the local-branch substitution from that of the peer fusion stage. All results report both detection performance and attained FPRs, and we run further ablation and fixed-cohort comparisons at different parameter budgets to understand the drivers behind these improvements. As detailed in Section~\ref{sec: protocol-history}, this protocol served as our testbed throughout development.
\subsection{Evaluation units and eligibility} 
The benchmark consists of 6 datasets partitioned into 7 strata: IoT-23~\cite{garcia2020iot23} (11 clients), TON\_IoT~\cite{moustafa2021ton} (9), Gotham~\cite{belarbi2025gothamdata,belarbi2026descriptor} (7), CTU-13~\cite{garcia2014ctu} (12), N-BaIoT~\cite{nbaiot2018,nbaiotdata2018} (9), and UNSW-MUD (IoT attacks) Group 1 (5) and 2 (5). Each of the resulting 58 evaluation targets correspond to some client that meets our experimental design criteria, and represent the set we wish to evaluate against (not e.g.\ an arbitrary draw of IoT clients). We specify precise subsets or complete dataset files used for each experiment, their sources and any adaptations made in Appendix~\ref{app:data}. Evaluation targets have different features sets and the associated peer clients will be drawn from the corresponding prepared set, i.e.\ no cross-dataset peer transfer is considered.

Each client dataset has an associated manifest file listing which of its targets we intend to evaluate, from which we draw the reported evaluation targets. These reported clients must have one or more attacks represented in their retained evaluation set, have at least 350 calibration samples and at least 4 eligible peer clients. To ensure sufficient evaluation attack rows, we discard a given attack if it falls short of a configurable threshold (default 600 samples), relaxed to 400 for IoT-23, TON\_IoT and CTU-13, and 30 for UNSW-MUD G1/G2. Hence a client's qualification depends on some knowledge of the evaluation attacks it will encounter (although it makes no use of this in fitting/calibration), which should be borne in mind when discussing the possibility of deployment with benign traffic only.

The 58 target clients have between 769 and 80k samples assigned to their benign train set and contain between 4 (UNSW) and 11 (CTU-13) peer clients. We note in particular for Gotham that there are 7 clients in each the peer and reported-sets, i.e.\ these are not necessarily one-to-one. All clients undergo specific pre-processing steps and row-subsampling prior to role-assignment. While rows follow their capture-time order (or as-given sequence for non-time-series sources), we do not directly consider detecting on the full contiguous stream of incoming observations.

\begin{table}[!htbp]
\centering
\small
\caption{Prepared target units and feature dimensions. Counts describe eligible reported targets; peer pools can be larger. Features and sampling are adapter-specific.}
\label{tab:corpora}
\begin{tabular}{llrr}
\toprule
Stratum & Evaluation unit & Targets & Features \\
\midrule
IoT-23 & Capture/scenario & 11 & 36 \\
TON\_IoT & Destination host & 9 & 29 \\
Gotham & Device capture & 7 & 37 \\
CTU-13 & Botnet scenario & 12 & 38 \\
N-BaIoT & Device & 9 & 115 \\
UNSW G1 & Device & 5 & 15 \\
UNSW G2 & Device & 5 & 13 \\
\bottomrule
\end{tabular}
\end{table}
 

Evaluation targets and the number of original input columns per corpus are summarised in Table~\ref{tab:corpora}. Note the TON\_IoT dataset also contains non-consumer IoT destinations such as servers/routers/general PCs. Further dataset subsetting/preprocessing details are provided per-corpus in Appendix~\ref{app:data}. As mentioned, any subsampling of samples occurs on retained rows.

The Gotham dataset in particular merits closer consideration. 5 input features aggregate over each successive full second of packets prior to retention (removing rows with no matching label or down-sampled via reservoir-sampling). Consequently, any benign retained sample could encapsulate information on subsequent attack/unlabelled rows within its time period, with rows after its time-point also being included in previous rows' features. Therefore these samples will span role-assignment splits and the underlying data is not exclusively from that role-period, however we do not make the claim that attack samples were used in training/calibration or that the protocol can capture causal detections at the timescale of successive packets.

\subsection{Role separation and fitting budgets} 
Samples for each prepared client are sorted by time/order as-given (deterministic) and we take the $\lfloor 0.4n\rfloor, \lfloor 0.6n\rfloor, \lfloor 0.8n\rfloor$ samples as boundaries between [train|reference|calibration|test] benign sets. A target's local channel(s) are trained on at most the first $N$ assigned train-set samples, and are removed from evaluation if fewer samples are available (no borrowing from later role sets is conducted).

Retained attack samples for each qualified attack-type are likewise deterministically sorted and the first and second half each contribute at most 10k samples to their peer-training and peer-evaluation sets respectively. In each case we exclude the evaluated target client from the corresponding set (to ensure no supervision from that client on attacks). The remaining peer set could still share the same attack labels as exist for that target; in practice only 11 of 258 target/attack-type pairs encountered had zero presence of that label in their training peer-set. Nonetheless, we make no stronger statement of label-disjointness nor of identical attack behaviour and instead refer to the protocol as detecting in the absence of attack samples at the target. Sample splits here occur sequentially and separately per client and label -- there is no unified global split.

We fit a pooled scaler using all loaded clients' train-sets prior to budgeted sample assignment to target channels. Samples assigned to target reference/calibration sets are kept separate from its test set. However, the total amount of target data involved in the protocol is therefore $>$ $N$ samples per Section~\ref{sec: problem}; this also extends to the peers due to shared scaling. We interpret the tested budgets of $N=100,400$ only in the context of fitting the local channel(s).

To limit any confounding of budget variation with changes to the evaluated target set, we conduct:
(i) the main comparison of $N=100$ vs.\ 400 on all 58 targets;
(ii) additional robustness checks of $N=750$ for the same set of 58 targets;
(iii) $N=100,400,2000$ for the subset of 47 targets containing $\geq 2000$ samples in their train set; and
(iv) $N=100,400,2000,20k$ for the subset of 15 targets containing $\geq 20k$ samples. All experiments compare paired budgets within each applicable subset.

\subsection{Configuration and comparison rules} 
Experiments are conducted for both alphas $\alpha \in \{0.01, 0.05\}$ with reserve $=0.5$, blocksize $=250$, max multiplier $= 64$ and bisection $= 20$. Training seeds $\{6,7,8\}$ were randomly chosen, with component-level seeds generated from hash of the target and relevant component. Varying these represents multiple realisations of learning with a fixed set of data-role allocations. We note that despite the method defaulting to one alpha value, we consider both values within this study campaign.

Our main alternative is a variant of PACE replacing its local branch(es) by an ensemble of 7 local methods whose mean vote is thresholded using all local parameters. As noted previously, it shares the same peer components, relative-evidence and target-peer calibration method, role allocation and seed assignments/number of targets under test. The local ensemble adds a 7$^{\text{th}}$ detector consisting of an Autoencoder (2 hidden layers, sizes 20 and 10, mirror output, 30 epochs, batch-size=256, lr=0.001). Comparisons to this configuration thus represent the total change of substituting PACE's local channel(s) for the anchor + bank.

We further report on two alternative configurations sharing the same local branch as PACE-alt. In these configurations, peer-probabilities are converted to an evidence w.r.t.\ the target reference probability, reduced to a maximum score and calibrated. Compared to the target-relative margin of Eq.\ref{eq: evidence}, this score instead offsets by $\ln(\max_{i}\frac{p_i}{ref_p}) + \ln(n+1)$. This 'max'-fusion configuration provides the reference point for peer score fusion, while a Bonferroni-style configuration similarly adjusts PACE's peer evidences using a uniform margin-evidence threshold at a fixed multiplier of 1. Both therefore share the use of peer evidence in their rule and should not be interpreted as removing the local channel from PACE.

The ablations on mean/max fusion of peer evidences replace PACE's peer rule with the given alternative, holding all other aspects of PACE constant. In addition, we therefore include PACE-alt's max-fusion baseline for comparison (which similarly replaces only PACE's peer rule but with a different conversion function and local channel configuration). Two further configurations assign all of the target's allocated nominal parameters either exclusively to the peer channel union or to the local branch alone ($7^{\text{th}}$ detector included). Note these do not match directly with the constituent terms within PACE. Finally, we explore dropping each individual channel (feature, rare and response); omitting the anchor channel; merging the anchor and bank as a single weighted union; fitting the response bank on evidences; and reducing PACE's reserve allocation to 0.25. All configurations use the appropriate role-allocated data and training seeds for their respective target/cohort/subset. Full description of all 14 configurations executed (incl.\ PACE, 12 internal study variations and the 2 comparison methods) including variants with unfavourable/null results can be found in Appendix~\ref{app:inventory}. They are intended as internal sanity-checks of our design choices; we have not reproduced any of the cited prior-art in our experimental setting to compare performance figures.

\subsection{Metrics and aggregation} 
Let $D(t,h,s)$ denote the proportion of attack test samples flagged for a given client $t$, attack $h$ and seed $s$. For each client and seed, the mean is calculated uniformly across all of its qualifying attacks, followed by a mean across seeds to yield the client-level detection. Client detections are further aggregated uniformly over all targets to form a pooled estimate (ensuring larger target datasets don't exert greater weighting; a consequence of a given stratum having more targets). FPR is defined similarly, as the proportion of test benign samples flagged per client/seed (average over seed), with an analogous pooling scheme across targets.

We define our primary statistic of interest as PACE's client detection minus that of PACE-alt (paired). We estimate a CI around the mean using $B=10k$ resamples (rng=314159) and CL=0.95, drawing (with-replacement) clients at the stratum-group level based on its unique ID ('enum', 2 UNSW groups are coalesced), with each client group's size and contribution reflected in its resample weight. Reported CI endpoints use the linear interpolation method on the bootstrap distribution of the mean. This ignores the effect of any shared peers/common-capture etc.\ but represents the target-level estimates across our chosen set of training seeds and attack test sets. There is no multiplicity correction for the remaining tests/analyses.

The alpha-multiplier sweep ranges over $\{0.55, 0.75, 1, 1.4\}$ and $\{0.4, 0.75, 1, 1.4\}$ for $\alpha=\{0.01, 0.05\}$ respectively, applied jointly to both of PACE's allocated parameter scales and the full PACE rule including any early termination is evaluated at each combination. The primary result adopts the rate-interpolation across sweep-points (capped at observed min/max) to best align with each target-seed's FPR. We compute summary quantiles of this target-seed mean FPR for PACE and compare against an exceedance measure of client-seed values $>2\alpha$.

\subsection{Internal protocol checks and reproducibility} 
\label{sec:protocol-history} 
Our pre-specified pass-fail check required: no clients to see a decrease of $>$ 0.05 in detection nor $>3\%$ of targets decreasing by $>0.01$; strictly positive lower CIs of the pooled improvement over PACE-alt for $N$ and swept budgets; increases to mean FPRs $\leq 0.004$ ($\alpha=0.01$) and $0.002$ ($\alpha=0.05$), increases in $>2\alpha$ client-seed FPR $\leq 0.03$, and passed similar summary checks removing any TON\_IoT targets, varying training seeds and each stratum in-turn. These criteria were intended as sanity-checks of the observed benchmark data rather than satisfying any external criterion or rigorous testing procedure.

A corresponding a-priori power study assuming normality around expected paired SD and an effect of interest of 0.02, with level/power = 0.05/0.8 was computed, though is not meant to represent any ex-post study of achieved power. E.g.\ it requires between 51-103 targets to detect the minimal effect for SD=0.05-0.10, which equates to detectable paired effects of 0.018393-0.036787 under our actual setting of 58 targets. Our analysis instead follows the observed paired differences and notes the caveats outlined above regarding its generalisability.

As noted, our development process used this evaluation set. One motivating example is a failure to detect on IoT-23 target (capture ID) 35-1 at an intermediate development stage which led to our anchor rule. Our project overview specifies the training seeds as uninformative (w.r.t.\ design) and our choice of evaluation seeds therefore provides only information on learn-variance stability when reused over an informative cohort. We do not assert these seeds provide any further statistical guarantees on the observed results.

Our first public commit to GitHub on 10 Oct 2026 contains the model implementation, all study variants, campaign specification and acceptance tests. Public tables and results were committed on 11th and a draft MS on the 12th. All UTC+01. We cannot substantiate the ordering of internal (private) commit history nor of results review with respect to these changes and hence cannot speak to any prospectivity of our experimental protocol.
